\documentclass[10pt,a4paper]{amsart}
\usepackage[margin=19mm]{geometry}
\usepackage{amsmath,amssymb,mathtools}
\usepackage[scr=rsfs,cal=euler]{mathalfa}
\usepackage{xfrac}
\usepackage[unicode,hidelinks,bookmarks=false]{hyperref}
\newcommand{\ie}{i.e.\ }
\newcommand{\cf}{cf.\ }
\newcommand{\resp}{respectively\ }
\newcommand{\Subsection}{\subsection}
\providecommand{\nobreakdash}{\nobreak\hbox{-}}
\setlength{\emergencystretch}{2em}
\multlinegap0pt
\usepackage{bm}
\usepackage{bbm}
\usepackage{dsfont}
\usepackage{xspace}
\usepackage{cancel}
\usepackage{mathtools}
\usepackage{amsthm}
\makeatletter
\@ifundefined{lemma}{\newtheorem{lemma}{Lemma}}{}
\makeatother
\DeclareMathOperator{\init}{Init}
\title[Semiring identities in the semigroup $B\_2$]{Semiring identities in the semigroup $B\_2$}
\author{Vyacheslav Yu. Shaprynski\v{i}}
\date{Source: arXiv:2604.01706v1 (CC BY 4.0)}
\begin{document}
\maketitle

\noindent\emph{Source and licence.} Semiring identities in the semigroup $B_2$. Version arXiv:2604.01706v1 (CC BY 4.0), distributed under the Creative Commons Attribution 4.0 International licence, as recorded in the arXiv licence metadata for this exact version and shown on the abstract page https://arxiv.org/abs/2604.01706v1. Full rights notice: licenses/arxiv-2604.01706-CC-BY-4.0.txt.\par\smallskip

\noindent\textit{Editorial scope.} Counted unit: the Lemma whose source label is \texttt{first letters} in the article (source0.tex lines 289--291, source label \texttt{first letters}), delivered together with the article's own complete original proof of it (source0.tex lines 292--334, one \texttt{proof} environment of 43 lines). No mathematics is written, repaired or re-derived: statement and proof are the article's own text line for line. Required context retained uncounted: crossing (lines 124--129); rook (lines 124--129); adding inequalities (lines 252--254); easy (lines 258--275); deduction (lines 261--263). Every definition, display and label of those lines is retained unchanged. Omitted from this package and not redistributed: 1--123, 130--251, 255--257, 264--288, 335--439. Nothing inside a retained excerpt is omitted or altered: every retained body is delivered exactly as the article's lines are written, so no omission receipt of any kind accompanies this package. Citations: the retained excerpts carry no citation command at all, so no bibliography entry of the article is transcribed and no reference is redistributed. Reference adaptation: none was needed and none was made; every reference inside the delivered unit resolves inside it, so no unresolved reference or citation is produced. Printed slips kept literal: no printed word, symbol, hypothesis or number of the counted statement or of its proof was altered. Layout and class: the article's own class is \texttt{amsart} with packages including amsmath, amssymb, amsthm, cite; the wrapper is the lane's pinned \texttt{amsart} editorial wrapper at 10pt on A4, which restates the theorem-like environments the retained text needs and adds the article's own macro definitions that the retained text needs (\texttt{\textbackslash init}), with extra wrapper packages (bm, bbm, dsfont, xspace, cancel, mathtools). The delivered numbering is the unit's own. Assets: no retained span contains a figure environment, a graphics-inclusion command, a PDF-page inclusion, a table or a TikZ picture; no image file is copied into this package, the package declares no asset dependency and no image file is redistributed with it. Licence: version arXiv:2604.01706v1 of this article is distributed under the Creative Commons Attribution 4.0 International licence, as recorded in the arXiv licence metadata for this exact version and shown on the abstract page https://arxiv.org/abs/2604.01706v1. A full rights notice is delivered at licenses/arxiv-2604.01706-CC-BY-4.0.txt. Characters: every retained excerpt is pure ASCII, so the delivered engine, XeLaTeX with its default Latin Modern text font, reports no missing character.
\begin{align}
\label{rook}
&x_2z_2\le x_1z_1+x_1z_2+x_2z_1,\\
\label{crossing}
&x_1yz_2\le x_1yz_1+x_2yz_2,
\end{align}

\begin{equation}
\label{adding inequalities}\mathbf p\le\mathbf r\text{ and }\mathbf q\le\mathbf  r\text{ imply }\mathbf p+\mathbf q\le\mathbf r.
\end{equation}

\begin{lemma}
\label{easy}
\textup{1)} If $((x,i),(y,j))\in\rho(\mathbf p)$ and $(x,i)\neq(y,j)$ then there exists a sequence
\begin{equation}
\label{deduction}(z_1,n_1),(z_2,n_2),\dots,(z_k,n_k)
\end{equation}
such that $(z_1,n_1)=(x,i)$, $(z_k,n_k)=(y,j)$, and
$$((z_l,n_l),(z_{l+1},n_{l+1}))\in\rho_1(\mathbf p)\cup\rho_2(\mathbf p)\cup\rho_3(\mathbf p)\cup\rho_3(\mathbf p)^{-1}$$
for $l=1,2,\dots,k-1$.

\textup{2)} For $l=1,2,\dots,k-1$ we have
\begin{align*}
&((z_l,n_l),(z_{l+1},n_{l+1}))\in\rho_1(\mathbf p)\iff n_l=1\text{ and }n_{l+1}=1,\\
&((z_l,n_l),(z_{l+1},n_{l+1}))\in\rho_2(\mathbf p)\iff n_l=2\text{ and }n_{l+1}=2,\\
&((z_l,n_l),(z_{l+1},n_{l+1}))\in\rho_3(\mathbf p)\iff n_l=2\text{ and }n_{l+1}=1,\\
&((z_l,n_l),(z_{l+1},n_{l+1}))\in\rho_3(\mathbf p)^{-1}\iff n_l=1\text{ and }n_{l+1}=2.
\end{align*}
\end{lemma}

\begin{lemma}
\label{first letters} Let $\mathbf p$ be a polynomial and $x\in c(\mathbf p)$. If $(x,1)\in\init(\mathbf p)$ then there exists a word $x\mathbf w\le_{\Sigma}\mathbf p$.
\end{lemma}

\begin{proof}
Let $y$ be the first letter of a word in $\mathbf p$. We have $(x,1)\in\init(\mathbf p)$ and $(y,1)\in\init(\mathbf p)$, so $((y,1),(x,1))\in\rho(\mathbf p)$. Using Lemma~\ref{easy}.1), consider the corresponding sequence~(\ref{deduction}). We use induction on the length $k$ of this sequence. 

\textbf{Induction base:} $k=1$. In this case $x=y$, so there exists a word $x\mathbf w$ in $\mathbf p$. In particular, $x\mathbf w\le_{\Sigma}\mathbf p$.

\textbf{Induction step.} Let $(z_t,n_t)$ be the last element of the sequence~(\ref{deduction}) such that $t\le k-1$ and $n_t=1$. There are three possible cases.

\emph{Case}~1: $t=k-1$. By Lemma~\ref{easy}.2), we have 
$$((z_{k-1},1),(x,1))\in\rho_1(\mathbf p).$$ 
By definition of $\rho_1(\mathbf p)$, there exists a word $x\mathbf w$ in $\mathbf p$.

\emph{Case}~2: $t=k-2$. The last three elements of the sequence~(\ref{deduction}) are $(z_{k-2},1),(z_{k-1},2),(x,1)$. By Lemma~\ref{easy}.2), we have 
$$((z_{k-2},1),(z_{k-1},2))\in\rho_3(\mathbf p)^{-1}$$ 
and 
$$((z_{k-1},2),(x,1))\in\rho_3(\mathbf p).$$
By definition of $\rho_3(\mathbf p)$, there exist words $\mathbf u_1z_{k-1}z_{k-2}\mathbf u_2$, $\mathbf v_1z_{k-1}x\mathbf v_2$ in $\mathbf p$. By induction hypothesis, we have $z_{k-2}\mathbf w'\le_{\Sigma}\mathbf p$ for some $\mathbf w'$. Hence
\begin{align*}
\mathbf p&\ge_{\Sigma}\mathbf u_1z_{k-1}z_{k-2}\mathbf u_2+\mathbf v_1z_{k-1}x\mathbf v_2+z_{k-2}\mathbf w'&&\text{by~(\ref{adding inequalities})}\\
&\ge_{\Sigma}\mathbf u_1z_{k-1}z_{k-2}\mathbf u_2+\mathbf u_1z_{k-1}x\mathbf v_2+z_{k-2}\mathbf w'&&\text{by~(\ref{crossing}) and~(\ref{adding inequalities})}\\
&\ge_{\Sigma}\mathbf u_1z_{k-1}z_{k-2}\mathbf u_2+\mathbf u_1z_{k-1}x\mathbf v_2+z_{k-2}\mathbf u_2&&\text{by~(\ref{crossing}) and~(\ref{adding inequalities})}\\
&\ge_{\Sigma}x\mathbf v_2.&&\text{by~(\ref{rook})}
\end{align*}

\emph{Case}~3: $t\le k-3$. Consider the elements $(z_t,1)$, $(z_{t+1},2)$, $(z_{k-1},2)$, $(x,1)$ of the sequence~(\ref{deduction}). By Lemma~\ref{easy}.2), we have
$$(z_{t+1},2)\,\rho_2(\mathbf p)\,(z_{t+2},2)\,\rho_2(\mathbf p)\cdots\rho_2(\mathbf p)\,(z_{k-1},2).$$
Therefore, there exist words $\mathbf w_1z_{t+1}$ and $\mathbf w_2z_{k-1}$ in $\mathbf p$. By induction hypothesis, there exists a word $z_t\mathbf w_3\le_{\Sigma}\mathbf p$. By Lemma~\ref{easy}.2), we have 
$$((z_t,1),(z_{t+1},2))\in\rho_3(\mathbf p)^{-1}$$ 
and 
$$((z_{k-1},2),(x,1))\in\rho_3(\mathbf p).$$ 
Therefore, there exist words $\mathbf u_1z_{t+1}z_t\mathbf u_2$ and $\mathbf v_1z_{k-1}x\mathbf v_2$ in $\mathbf p$. We have
\begin{align*}
\mathbf p&\ge_{\Sigma}\mathbf w_1z_{t+1}+\mathbf u_1z_{t+1}z_t\mathbf u_2+\mathbf w_2z_{k-1}&&\text{by~(\ref{adding inequalities})}\\
&\ge_{\Sigma}\mathbf u_1z_{t+1}+\mathbf u_1z_{t+1}z_t\mathbf u_2+\mathbf w_2z_{k-1}&&\text{by~(\ref{crossing}) and~(\ref{adding inequalities})}\\
&\ge_{\Sigma}\mathbf w_2z_{k-1}z_t\mathbf u_2,&&\text{by~(\ref{rook})}
\end{align*}
whence
\begin{align*}
\mathbf p&\ge_{\Sigma}\mathbf w_2z_{k-1}z_t\mathbf u_2+\mathbf v_1z_{k-1}x\mathbf v_2+z_t\mathbf w_3&&\text{by~(\ref{adding inequalities})}\\
&\ge_{\Sigma}\mathbf v_1z_{k-1}z_t\mathbf u_2+\mathbf v_1z_{k-1}x\mathbf v_2+z_t\mathbf w_3&&\text{by~(\ref{crossing}) and~(\ref{adding inequalities})}\\
&\ge_{\Sigma}\mathbf v_1z_{k-1}z_t\mathbf u_2+\mathbf v_1z_{k-1}x\mathbf v_2+z_t\mathbf u_2&&\text{by~(\ref{crossing}) and~(\ref{adding inequalities})}\\
&\ge_{\Sigma}x\mathbf v_2.&&\text{by~(\ref{rook})}
\end{align*}
\end{proof}

\end{document}
